% ================================================================
%  第 8 章 师与生 (原书 p.167-174, 边码 p.57-p.66)
%  源文: _work/src1/task-ch08a.md
% ================================================================
\chapter{师与生}

\section{(25) 学生与研究计划}

% ---------------- p.167 ----------------
\origpage{57}我还没有把当年与其他数学家的关系巡视完——那时我觉得自己与他们同属一个世界,同一个“数学共同体”。剩下要看的,主要是我与我那些学生的关系——按我所经历的样子——以及我与另一些人的关系,那些人把我看作前辈。

总的来说,我相信可以毫不保留地说:我与学生的关系,是尊重的关系。至少在这一件事上,我相信,我当年做学生时从前辈们那里得到的东西,并没有随着岁月而变坏。由于我素有做“难”数学的名声(说实话,这是个再主观不过的概念!),而且被人认为比别家老板更苛刻(这一点就主观得少些了),来找我的学生从一开始就相当有劲头:“他们是想干的”!只有过一个学生,开头有点“提不起劲”,说不准他能不能起步——后来还是起了,而且没等我推一把,他自己就发动起来了。

就我所能记起的而言,凡是要求跟我工作的学生,我都收下了。其中有两个,几周或几个月之后发现我的工作方式不合他们的脾胃。说实话,现在我觉得那两次都是“卡住”的局面,而我当时却草率地把它读成了不适于数学工作的迹象。今天,我要作这类预言会谨慎得多。我毫不犹豫地把自己的印象告诉了这两位当事人,劝他们不要在一个——据我看——不合他们天性的行当里继续下去。事实上我后来得知,至少这两个学生中的一个,我是弄错了:这位年轻研究者后来在代数几何与数论交界处的一些难题目上赢得了名声。另一个学生是位年轻女子,她跟我这段不如意之后有没有继续,我就不清楚了。不能排除:我对她能力的看法说得太斩钉截铁,使她泄了气——而她也许和别人一样有能力做出好工作。我觉得我对这两个学生和对别人一样,是给了信任的。我缺的倒是%
% ---------------- p.168 ----------------
辨别力:面对那些肯定是“卡住”、而非“无能”的迹象时,没能分清两者\fnoteP{1}{(18)}{我想,这份缺乏辨别力,在这两回并不是出于我的疏忽,而更是出于不成熟、出于无知。直到大约十年之后,我才开始留意“卡住”的机制——在我自己身上,也在我亲近的人身上,在学生身上——并估量出它们在每个人的生活里扮演的巨大角色,而且不只是在中学或大学里。当然,我后悔这两回没有更成熟一些的辨别力;但我不后悔把自己那些印象——无论有没有根据——明明白白地说出来。当我在某一场合看出工作是敷衍着做的,把这些事照它们本来的样子点明,在我看来是必要而有益的。而如果在另一场合,我由此得出的结论是草率的、没有根据的,那么负有责任的也不只我一个。被这样摇了一下的学生,仍然有得选:要么从中吸取教训(也许头一回就是这样),要么就此泄气,也许还改了行(这也未必是坏事!)} (18)。

从六十年代初起,也就是在大约十年里,有十一个学生在我这里做了国家博士论文\fnoteP{2}{(19)}{\emph{耶稣与十二使徒}\par 从1970年至今,还有一个学生伊夫·拉德加耶里(Yves Ladegaillerie)在我这里准备并通过了论文。前一时期的学生是:P.贝尔特洛(P. Berthelot)、M.德马祖尔(M. Demazure)、J.吉罗(J. Giraud)、M.阿基姆夫人(Mme M. Hakim)、黄春生夫人(Mme Hoang Xuan Sinh)、L.吕西耶(L. Illusie)、P.茹阿诺卢(P. Jouanolou)、M.雷诺(M. Raynaud)、M.雷诺夫人(Mme M. Raynaud)、N.萨维德拉(N. Saavedra)、J.L.维尔迪耶(J.L. Verdier)。(其中六位还是在1970年之后才完成论文工作的,也就是在我很少能抽出时间指导他们的时期。)在这些学生当中,米歇尔·雷诺(Michel Raynaud)占一个特殊的位置:他论文工作所处理的那些本质性的问题与概念,是他自己找到的,而且他完全是独立地把它们发展出来的;所以我这个严格意义上的“论文导师”的角色,只限于读已完成的论文、组建评审委员会,并做其中一员。\par 在我自己出题目的时候,我很当心地只限于那些我足够熟悉、与我有足够深联系的题目,以便在需要时觉得自己能够支撑学生的工作。一个显著的例外是米谢勒·雷诺夫人(Michèle Raynaud)关于基本群的局部与整体Lefschetz定理的工作,这些定理是以适当的平展site上的1-栈的形式表述的。这个问题在我看来是困难的(后来也确实证明如此),而对于我提出的那些猜想(它们其实几乎无可置疑),我并没有证明的思路。这项工作一直延续到七十年代初;雷诺夫人(如同此前她丈夫的情形)发展出一套精细而独到的方法,没有得到我任何帮助,也没有得到别处的帮助。这项出色的工作还引出了一个问题:把雷诺夫人的结果推广到n-栈的情形——在我看来,在概形的框架里,这应当代表“弱Lefschetz定理”一类定理的自然归宿。不过,这里相应的猜想(它同样几乎无可置疑)的表述,本质上要用到n-栈的概念;而n-栈的追寻,据说正是本书的主要对象[其实那是《数学反思》第3卷,不是眼前这部《收获与播种》第1卷——见导论,p.(v)],正如“追寻栈”这个书名所表明的。我们大概会在适当的地方再回到这一点上,\par 另一个相当特殊的情形是黄春生夫人:我最初是1967年12月在河内遇见她的,那是我在被疏散的河内大学讲一个月“课程—讨论班”的时候。第二年,我向她提出了她的论文题目。她是在战争时期格外艰难的条件下工作的,与我的接触只限于零星的通信。1974/75年她得以到法国来(借温哥华国际数学家大会之机),并在巴黎通过了论文(评审委员会由嘉当(Henri Cartan)主持,成员还有施瓦茨(Laurent Schwartz)、德尼(J. Deny)、齐斯曼(M. Zisman)和我)。\par 最后我还得提到皮埃尔·德利涅和卡洛斯·孔图-卡雷尔(Carlos Contou-Carrère):这两位都多少算是学生,前者在1965–68年前后,后者在1974–76年前后。两人显然都有(而且一直有)不寻常的才能,他们运用这些才能的方式大不相同,运气也大不相同。来比雷之前,德利涅多少做过蒂茨(J. Tits)的学生(在比利时)——不过我怀疑,按通常意义上的“学生”来说,他在数学上不曾是谁的学生。孔图-卡雷尔做过圣塔洛(L. Santaló)的学生(在阿根廷),还曾有一段时间——多少算是——做托姆(René Thom)的学生!联系建立起来的时候,两人都已经有一个数学家的身量了;只是孔图-卡雷尔还缺方法、缺手艺。\par 我在数学上对德利涅的作用,只限于零零碎碎地让他知道我关于代数几何懂得的那一点点;而他学它,就像听人讲一个故事——仿佛他向来就知道。此外,还在于沿途提出一些问题,而这些问题他多半当场或在随后几天就找到了答案。这就是我所知道的德利涅最早的工作。1970年以后的(他的,以及我那些“正式学生”的),我仅凭极为零散、遥远的回声才有所知[其实那是《数学反思》第3卷,不是眼前这部《收获与播种》第1卷——见导论,p.(v)]。\par 按孔图-卡雷尔自己在论文开头所说的,我在他身边的作用,只限于把他引入概形的语言。无论如何,他近些年作为国家博士论文所准备的工作,我只是从很远的地方跟着;那题目是最时髦的那一类,超出了我的能力。在广阔世界里经历了几番波折之后,孔图-卡雷尔最近终于在最后一刻——而且(现在在我看来)是违着他的心意的——被逼得来求助于我,请我充当论文导师并组建评审委员会。(这就使他冒着一种风险:显得像是“1970年之后”的格罗滕迪克的学生——在这种情势下,这可能带来严重的不便……)我尽我所能履行了这项任务;这大概是我最后一次行使这个职能(在国家博士论文这一层面上)。在这个有点特殊的情形里,我格外高兴有让·吉罗(Jean Giraud)的友情协助:他也从自己的时间里抽出一两个月,细读那部篇幅庞大的手稿,并写了一份详实而热忱的报告。} (19)。各自选定了合意的题目之后,他们都干得起劲,而且(我是这样感觉的)都与自己选定的题目深深认同。然而有一个例外:有个学生选了一个“该做的”题目——也许并无真正的信念——而这题目也有它不讨人喜欢的一面:它要做的是把一批已经到手的思想作技术上的定稿整理,有时艰难,甚至干枯,再也没有多少意外、多少悬念可言,%
% ---------------- p.169 ----------------
\origpage{58}只是把它们摆出一个全景来\fnoteP{3}{(20)}{这让我想到莫妮克·阿基姆(Monique Hakim)选的题目:说实话,那也不更招人喜欢;我真不知道她是怎么保持住精神的!她固然有时也吃力,但无论如何没到让她变得忧郁或阴沉的地步,而我们之间的工作是在一种亲切、放松的气氛里做的。} (20)。

我当时被一项庞大计划的需要卷着走——我需要人手——所以在出这个题目时,必然是缺乏心理上的辨别力:它肯定不合这位学生特有的性情。他那一头呢,大概也没太看清自己上的是一条什么船!总之,他和我都没能及时看出:这事一开始就起步错了,不如换个题目重新来过。

显然,他干活没有真正的信念,而且始终带着一副多少有些忧郁、阴沉的神气,不曾放下。我想我当时已经到了这样一个地步:不太留意这类事情——而这类事情(我本该记得)在任何研究工作中,而且不只在研究工作中,都能造成天壤之别!我那时的角色,仅限于:工作显出拖沓的样子,我就心烦;工作又动了,我就宽慰地“呼!”一声;等到最后,预定的计划总算“收了尾”。

直到1970年我那次“醒悟”之后好几年,因为要和这位从前的学生通信(他当了教授——再说,在那宽厚的年月里,人人都当上了教授!),我才想到:那一回终究是有什么地方不对头,那也许算不上一次彻底的成功。今天,它在我看来是一次失败——尽管“计划收了尾”(绝非草草了事!),尽管有文凭,还有随之到手的位置。而且我要负很大一部分责任:我把一个计划的需要,放在了一个人的需要之前——而这个人,是怀着信任把自己交到我手上的。我刚才还拿来标榜的那种“尊重”(“毫无保留”!),我说我对学生们做到了——在这里始终是表面的,与尊重的真正灵魂相隔绝:那灵魂就是对一个人的需要怀着情意的关注,至少当这些需要的满足取决于我的时候。这里的需要,是工作中的一份快乐;没有它,工作就失了意义,变成强制。

在这次反思中,我有机会谈到过一个“没有爱的世界”,并且在自己身上寻找那个我拒斥的世界的胚芽。这里就有一粒,而且个头不小——至于它在别人身上怎样长起来,我今天就说不上来了。这种表面的、缺乏关注、缺乏真正爱的尊重,也是我给予自己孩子们的“尊重”。和他们在一起,我有这份特权:看着这粒种子长起来,看着它蔓延。我也多少明白了:面对收成,皱眉头是没有用的……

\section{(26) 严格与严格}

如果把这位学生除外——他肯定并不比别人“天分”差——我可以说,我与学生之间的关系是融洽的,常常甚至是亲热的。出于无奈,他们%
% ---------------- p.170 ----------------
\origpage{60}都学会了耐心对待我作为“老板”的两个主要毛病:一是字迹无法辨认(不过我相信,最后人人都学会了辨认我的字);二是——这一点当然更严重,而我很晚才发现——我从根子上难以跟随别人的思路,除非先把它译成我自己的意象,用我自己的风格重想一遍。我远更倾向于把自己深受浸润的某种对事物的眼光传给学生,而不是鼓励他们心中绽出一种个人的眼光——那眼光也许与我的颇为不同。与学生的关系中这层困难,到今天也还没有消失;但我觉得它的影响减弱了,因为我意识到了自己身上的这种倾向。也许我的气质——无论是先天的还是后天的——使我更适于孤独的工作(我起初十五年的数学活动,约1945至1960年,也正是这样过来的),而不那么适于做“师傅”,去和那些数学上的志业与个性尚未完全成形\fnoteP{4}{(21)}{也许更准确的说法是:就我这种气质而言,我至今缺的是必要的成熟,还不能完全承担起一个教育者的角色。我后天形成的气质,长久以来被一种过分占优势的“男性”(或“阳”)特征所标记;而成熟的一个方面,恰恰是以“女性”(或“阴”)为主的“阴阳”平衡。\par(后来补写。)比我缺一份成熟更甚,我看在我作为教育者的生活里,直到今天,我缺的是一种慷慨——一种比在时间和精力上的投入更为精细地表达出来的慷慨,也是更本质的慷慨。在我教学的第一时期,这份欠缺并没有以可见的方式表现出来(比如说,以一连串失败局面的积累),多半尤其是因为,那些选择来跟我工作的学生身上的强烈动机把它补偿了。相反,在第二个时期,从1970年至今,在我看来,对于我在研究层面(也就是从DEA这一层往上)的教学中所看到的整体失败,这份欠缺至少是原因之一,而且无论如何是最直接牵涉到我本人的那一项原因。关于这一点,见《纲领概要》第8段,以及第9段《一项教学活动的总结》,那里透出这项活动七八年来留在我身上的受挫感[并请对照后来补写的注(23iv)。]。} (21)的学生相处。不过也确实:我从小就喜欢教书;而且从六十年代到今天,我能有的那些学生在我生活里占了重要的位置。这也就是说,我的教学活动、我作为教师的作用,在我的生活里曾经有过、并且仍然保有一个很大的位置\fnoteP{5}{(22)}{不过也许不会太久了,因为我已经决定申请进入法国国家科学研究中心(CNRS),从而结束在大学环境里的教学活动——这几年来,那活动变得越来越成问题。} (22)。

在我教学活动的前一时期,我和任何一个学生之间都没有过表面上的冲突,哪怕只是关系里一阵短暂的“冷淡”也没有过。只有一次,我不得不对一个学生说:他工作不认真;要是照这样下去,我就没兴趣再和他一起干了。那事情究竟怎样,他当然和我一样清楚;他改了过来,这件事就了结了,没有留下一片乌云。另一次,已经是七十年代初,那时我最主要的精力投在“生存与生活”(Survivre et Vivre)小组的活动里;有个学生看了我刚写好的、关于他工作的论文评审报告(我习惯如此),发起火来,认为报告里某些议论是在质疑他工作的质量(这绝不是我的本意)。这一回是我把话收回、把方向拨正了,没有为难。当时我并不觉得这段小插曲会在我们的关系里留下阴影;不过也可能是我弄错了。这位学生和我的关系,比我和别的学生更少个人色彩(我说过的那个“忧伤的学生”除外);那只是一份良好的工作关系,仅此而已,我们之间并没有真正的暖意流过。不过,我并不认为是我身上一种无意识的不厚道,使我把那些他认为于己不利的议论写进了报告——他还说“这事他不会放过”,不像他的一个同伴,那人已经在我这里通过了论文。与另一位学生——他天性敏感而多情——我则有着一份%
% ---------------- p.171 ----------------
\origpage{61}格外亲热的关系;如果说我在关于他论文的报告里也写进了那类曾让他同伴大为不快的议论,那肯定不是出于不厚道!再说,对他们两位,也对我所有的学生,只要我对他们提交的工作不是完全满意,我就不会给答辩开绿灯。而且,那个时期我的学生,论文一通过,没有哪一个不是很快就找到了一个与他相称的位置。

直到1970年:我几乎随时都能抽出时间给学生\fnoteP{6}{(22')}{即便在1970年之后,当我对数学的兴趣在我的生活里变得零散而边缘时,我也不认为有过学生来找我、要我跟他一起工作而我推辞掉的场合。我甚至可以说,除了两三个例外,1970年之后我的学生对他们所做工作的兴趣,远低于我自己对他们题目的兴趣——即便在我除了去学院的日子几乎不操心数学的那些时期。因此,我在1970年之前随时能抽出时间给学生的做法,以及作为它主要标志的那种在工作上的极端严格,对于我后来大多数学生都会毫无意义——他们做数学没有信念,仿佛是靠着一种他们不得不对自己不断施加的努力……} (22')。时机成熟的时候,每逢用得着,我就和这一个或那一个待上一整天、必要时好几天,去啃那些还没弄妥的问题,或者一起把他们文稿的一稿又一稿过一遍。就我对这些工作时段所体验到的而言,我不觉得自己在里面扮演过拍板定案的“导师”角色;每一次都是共同的研究,讨论平起平坐地进行,直到双方都完全满意。学生投入的精力是巨大的,当然无法与我自己应投入的相比——不过我这边有更多的经验,有时还有训练得更敏锐的嗅觉。

然而,在我看来,对任何研究——无论智性的还是别的——的质量最要紧的东西,绝不是经验的问题。那是对自己的严格要求。我要说的这种严格,本质上是精细的;它不属于一丝不苟地遵从无论什么样的规范——严格的规范也好,别的规范也好。它在于对我们自身内部某种细微之物的极度关注,而那东西逃开一切规范、一切尺度。那细微之物,就是:对所考察之物的理解,有,还是没有。更确切地说,我要说的这种关注,是对每一时刻当下那种理解之品质的关注:从一堆杂乱堆叠的概念与陈述(假设的或已知的)发出的那片噪音,直到完全的满足、完美理解那种完成了的和谐。一项研究的深度——无论它的结果是片断的理解还是总体的理解——就在这份关注的品质里。这样一种关注,并不像遵守某条训诫、凭着“当心点”的刻意意图去留神的结果;在我看来,它是从求知的激情里自发地生出来的,是那些把求知的冲动与它那些自我中心的赝品区分开来的标志之一。这种关注有时也被叫作“严格”。这是一种内在的严格,独立于某一时刻在某门(不妨说)特定学科里可能通行的那些严格规范。如果说在这本书里我容许自己对那些严格规范(我曾教过它们,它们自有其存在的理由和用处)放肆一些,我并不认为这本书里这种更本质的严格,比我过去那些规范文体的著作里更少。而如果说我毕竟也许还能传给学生某种比一套语言和一种手艺更贵重的东西,那无疑就是这份严格、这份关注、这份严谨——如果说不是在与他人的关系和对自己的关系里(在那一层,它在我身上和在任何别人身上一样缺),那至少是在数学工作里\fnoteP{7}{(23)}{\emph{孩子与老师}\par“传递”这个词在这里并不真正符合事情的实际,它提醒我要采取一种更谦逊的态度。这种严格不是可以传递的东西,至多只能被唤醒或鼓励——而它从很小的年纪起就被忽视了、被挫伤了,被家庭环境,也被学校和大学。就我所能记起的,这种严格一直在我那些追寻里——至少是智性性质的追寻里;而我不认为它是我的父母传给我的,更不认为是学校的老师们、或我那些数学前辈们传给我的。在我看来,它属于“无辜”的禀性,因而也属于每个人生来就被分派到的东西。这份无辜很早就“青的、没熟的都尝遍了”,于是它不得不或深或浅地沉下去,以致在余生里往往再也见不到它多少痕迹。在我身上,由于一些我还没想到去探究的原因,某种无辜在智性好奇这个相对无害的层面上存活了下来;而在别处一切地方、在所有人身上,它都深深地沉了下去,神不知鬼不觉!也许“教育”——就其完整意义而言——的秘密,或者说谜,就在于重新与这份表面上消失了的无辜取得接触。但是,如果在教育者本人身上这份接触不是已经先在、或已被重新找到,那就谈不上在学生身上重新找到它。而这时教育者“传”给学生的,绝不是这份严格或这份无辜(它们在这两人身上都是天生的),而是一种尊重,一种对那件通常被弃置之物的默默的重新肯定。} (23)。

% ---------------- p.172 ----------------
\origpage{62}这确实是一件很微薄的东西,但也许毕竟聊胜于无。

\section{(27) 失误——或二十年后}

也许除了我谈过的那两个学生——与他们终究没能建立起工作关系——之外,我不记得其他那些来找我、要求跟我工作的学生,是带着“怯场”或畏惧来的。想必他们多少已经认得我,哪怕只在我于法国高等科学研究所(IHÉS)的讨论班上待过一阵子。倘若在我们关系开始时有什么拘束,它最后总会在工作过程中消散,不再留下痕迹。不过这里我要举两个例外。一个就是那位始终没能真正对工作生出兴味的学生,甚至在我们共同工作时,他也一直只是单音节地应声。也许他来的时候,我能留给学生的时间正要减少;和他也没有过那种对着稿子、一坐就是一个下午、一整天的时段。不,确实,我不记得有过那样的时段;我倒觉得我们多半是匆匆见一面,一两个钟头,看看他做到哪儿了。说到底,撞上我,他大概是最倒霉的一个!

我本来要说的另一个学生,则是在我还随时能抽出时间给学生的那个时期跟我工作的。我们的关系从一开始就融洽。他甚至属于那少数几个与我建立起友情的学生之列——我有时去他们家,他们也来我家,有点像两家之间的往来。确实,即便在这些情形里,关系也始终停留在相对表面的层次——至少就我这一方而言。在意识层面上,那时我对自己家里、自己屋顶下发生的事就已经所知无几;至于我那些数学家朋友——学生也好,不是学生也好——的生活,除了妻子和孩子的名字,我最终几乎一无所知(就连这些名字,我有时也会忘掉,而从没有人为此怪我!)。也许我曾是“书呆子”的一个极端例子;但我认为,在我所认识的那个数学圈子里,大多数关系——即便是友好的、亲切的——都停留在这种表面的层次:彼此最终所知极少,除了在“未言明”的层面上被感知到的那一点。这无疑是人与人之间的冲突在这个圈子里如此罕见的原因之一——尽管在我看来很清楚:分裂存在于我的大多数同事和朋友内部,也存在于他们的家庭内部,和我身上、和别处一样多。

我不认为我与这个学生的关系有别于我与别人的关系;当时我也不觉得反过来,他与我的关系与其他学生的关系有什么显著的差别,尤其是与那些同我结下友情的学生相比。只是不久以来,我才得以明白:那一定是一份比我对大多数别的学生更强的关系。一个未曾说出的冲突,其可见的表现像一次意外的启示般到来——那已是他做我学生将近二十年之后。直到那时,我才把它和一件早已被遗忘的“小”事联系起来。在很长一段时间里,也许甚至在整个那一段时期(也就是几年)里——%
% ---------------- p.173 ----------------
\origpage{63}\origpage{64}我们或多或少定期一起工作的那些日子——这个学生一直保留着某种“怯场”。它每次见面都表现出来,带着不会骗人的迹象。随后,在共同工作的过程中,这些迹象又相当快地消失了。这些不安的迹象当然让我不自在,而我觉得他更不自在。我们两个都装作不知道这事,这也理所当然。想必我们谁也不会想到去谈它,甚至不会想到在心里给这样一个古怪、显然值得一顾的处境一点注意!在他那边,也在我这边,这份“怯场”想必被感受为一处单纯的“失误”,一处本不该有的东西。那“失误”定期提醒我们想起它,但每一次,它都识趣地消失,让我们有时间安安稳稳地去忙正经事——数学——同时忘掉“那本不该有的东西”。我不记得自己哪怕一次停下来,对那失误的意义提出什么问题;我也确信,我的学生兼朋友那一头也是如此。想必在我们各自从幼年起在身边所见的一切里,没有任何东西能在他或我心里提示出另一种对待一件令人不快之事的态度,除了尽可能把它挪开、让它不再碍事。在那件事上,这完全做得到,甚至很容易;而我们完全一致地什么都没看见、什么都没感觉、什么都没听见。

然而,借这两三年来回到我这里的许多回声与相互印证,我明白:那被我们当作“本不该有”而挪开的东西,并不因此就不再存在、不再表现出来。有时回到我这里的东西,也同样“不该有”——可它“在”,而且现在已经不能一手背就把它拂开了……

\section{(28) 未完成的收获}

直到1970年第一次“醒悟”之前,我与学生的关系,正如我与自己工作的关系,是满足与快乐的源泉,是我生活中那份和谐感的可触、无可否认的基础之一——那份和谐继续赋予我的生活以意义,而与此同时,一种捉摸不定的毁坏正在我的家庭生活里肆虐。那时候,在我眼里这些关系里没有任何表面冲突的成分;当时它们当中没有哪一个,哪怕在一瞬之间,成为挫折或痛苦的起因。有一件事也许显得悖理:我与某个学生的关系里的冲突,偏偏要等到那次著名的醒悟之后才显出来——也就是在一个转折之后,那转折给了我的生活一种它从前不曾有过的开阔,也许还给了我这个人一点点柔韧的开端——人们会想,这些品质本该有助于化解或避免冲突,而不是挑起或激化它。

可是再仔细看,我就看出这悖理只是表面的;无论从哪个角度看,它都消散了。第一点:一个冲突要有机会化解,首先得它已经表现出来。冲突表现出来的阶段,相对于隐藏着、被忽视的冲突而言,是一种成熟;而后者其实也有表现,而且那些表现恰恰因为借它们表达出来的冲突仍被忽视而更为“有效”。第二点:一个冲突要能以可辨认的方式表现出来,首先得有一段距离缩小了,或者消失了。近十五年来我生活里发生的种种变化——尤其是在一次次“醒悟”过程中——在我看来,都是那种能缩小距离、抹去隔绝的变化。一个冲突面对一位有声望、受人仰慕的老板时难以表达,面对一个被剥去了权力位置的人(在这里是自愿剥去的)、一个从某个握有权威与声望的圈子里自我放逐的人,就自在多了——这个人越来越不被看作某种实体(比如数学)的化身或特权代表,而越来越被看作一个和别人一样的人:%
% ---------------- p.174 ----------------
\origpage{65}\origpage{66}一个不仅可能被伤到,而且越来越不愿把自己的伤或痛藏起来的人。第三点,也是最要紧的:我第一次醒悟以来的演变——尤其是在那个时期以及随后的几年里——本就容易在我那些旧学生井井有条的世界里激起(或者也许是唤醒)疑问、不安、一种“重新审视”。我有充分的机会看出:不仅在他们身上如此,在数学界我昔日的朋友与同伴中间也如此,有时甚至在一些只凭传闻认识我的科学同行中间也如此。

还得说:一个稍有深度的冲突能得到化解,是最罕见不过的事。最常见的情形是:尽管有一切休战与表面的和解,我们那些冲突越排越长的队伍仍一路跟着我们,几乎寸步不离地跟一辈子,最后才在殡葬工那双阴沉的手里把我们松开。我有时有幸看到一个冲突多少解开了,有时甚至看到它在认识中化解——但直到如今,这样的事还不曾在我与某个学生、或与数学界某位昔日朋友的关系中、并借这关系而发生。我也很清楚:这样的事是否终会发生,毫无把握,哪怕我还能再活一百年。

有一件事很值得注意:我与某段过去决裂的那个时刻——我是说离开IHÉS那个插曲(这个机构多少像是我周围形成的那个数学小宇宙的“母模”)——这个决定性的插曲,同时竟是我某个学生对我的深刻敌意第一次表现出来的场合。肯定是这一情形,使那个插曲格外难堪、格外痛苦,像一次分娩或一次出生发生在格外艰难的条件下。当然,那时我还不能在我后来学会的光亮里去看这个插曲——它的意义当时逃开了我。此后很久,这份痛苦的意外仍留在那里。然而,就在那一年的夏天,这场带着苦涩的离开已经显为一次解放——就像一扇门忽然大大敞开(只消我推它一下!),通向一个不曾料想的世界,唤我去发现它。此后每一次新的醒悟,也都是一次新的解放:发现一种从属、一种内在的束缚,并且重新发现一个巨大未知者的在场——它就藏在那本该“已知”之物的熟悉外表后面。但同样在这整整十五年里,直到今天,这份固执、沉默而无懈可击的敌意一直跟着我,成为我在数学生活里所知道的唯一而持久的挫折之源\fnoteP{8}{(23')}{不过七八年来,在我的数学生活里还有另一个慢性的“受挫之源”,只是它在这些年里表达得隐蔽得多。靠着同一种“令人受挫”的局面在我的教学活动中顽固地重复、积累,它最终还是显了出来,并且最后爆发成一种“受够了!”,使我几乎终止了一切所谓“研究指导”的活动。在这次反思的过程中,我有一两次触及这个问题,最后至少在末尾稍作考察。我在那里至少描述了这份受挫感,并考察了它在我“回到数学”这件事上扮演的角色(参看第50段“一段过去的重量”)。} (23')。也许我可以说,它是我为那第一次解放、以及随后的几次解放付出的代价。但我很清楚:解放与内在的成熟,是与“要付的代价”不相干的事;它们不是“赚”与“亏”的问题。或者换个说法:当收获被带到头、当它完成的时候,就没有损失——那看似“损失”的,已经变成了“收获”。而事情就变得清楚:我至今还没能把那场收获带到头;就在我写下这几行的此刻,它仍是未完成的。
